% ============================================================
% PARTE VIII - ARMONIA CROMATICA
% CAPITOLO 38 - LA SESTA NAPOLETANA
% ============================================================
\part{Armonia cromatica}

\chapter{La sesta napoletana}
\label{chap:sesta_napoletana}

La \textbf{sesta napoletana} è il primo degli accordi cromatici di preparazione alla dominante. Non appartiene alla scala diatonica: si ottiene abbassando di un semitono il secondo grado e costruendovi sopra una triade maggiore, che si usa quasi sempre in primo rivolto (da cui il nome di «sesta»).

\section{Struttura}
\label{sec:struttura_napoletana}

\begin{definizione}[Sesta napoletana]
	\label{def:sesta_napoletana}
	La \textbf{sesta napoletana} (\sigla{II$\flat$}{6}{}) è la \textbf{triade maggiore costruita sul II grado abbassato} della scala, in \textbf{primo rivolto}: al basso sta la terza, cioè il IV grado. In La minore: Si$\flat$-Re-Fa con il Re al basso. In Do maggiore (accordo preso a prestito dal modo minore): Re$\flat$-Fa-La$\flat$ con il Fa al basso.
\end{definizione}

Il II grado abbassato si chiama \textbf{nota napoletana}. L'accordo ha funzione di \textbf{sottodominante} (sostituisce il \sigla{II}{6}{} o il IV) e conduce alla dominante, direttamente oppure attraverso il 6/4 di cadenza (Capitolo \ref{chap:quarta_sesta}).

\begin{esempio}[Costruzione della sesta napoletana]
	\label{ex:costruzione_napoletana}
	La triade maggiore sul II$\flat$ in stato fondamentale e in primo rivolto, in Do maggiore e in La minore.
	\musica{es_cap21_01}
\end{esempio}

La sesta napoletana vale sia per la tonalità maggiore sia per la minore omonima (Do maggiore e Do minore hanno entrambe \sigla{II$\flat$}{6}{} $=$ Fa-La$\flat$-Re$\flat$). Nelle tonalità con molte alterazioni la grafia corretta richiede doppi bemolli.

\begin{center}
\small
\begin{tabularx}{\linewidth}{lLLLL}
\toprule
\textbf{Tonalità} & \textbf{II$\flat$ (fondamentale)} & \textbf{IV (al basso)} & \textbf{VI} & \textbf{\sigla{II$\flat$}{6}{} (dal basso)} \\
\midrule
Do & Re$\flat$ & Fa & La$\flat$ & Fa--La$\flat$--Re$\flat$ \\
Sol & La$\flat$ & Do & Mi$\flat$ & Do--Mi$\flat$--La$\flat$ \\
Re & Mi$\flat$ & Sol & Si$\flat$ & Sol--Si$\flat$--Mi$\flat$ \\
La & Si$\flat$ & Re & Fa & Re--Fa--Si$\flat$ \\
Mi & Fa & La & Do & La--Do--Fa \\
Si & Do & Mi & Sol & Mi--Sol--Do \\
Fa$\sharp$ & Sol & Si & Re & Si--Re--Sol \\
Re$\flat$ & Mi$\flat\flat$ & Sol$\flat$ & Si$\flat\flat$ & Sol$\flat$--Si$\flat\flat$--Mi$\flat\flat$ \\
La$\flat$ & Si$\flat\flat$ & Re$\flat$ & Fa$\flat$ & Re$\flat$--Fa$\flat$--Si$\flat\flat$ \\
Mi$\flat$ & Fa$\flat$ & La$\flat$ & Do$\flat$ & La$\flat$--Do$\flat$--Fa$\flat$ \\
Si$\flat$ & Do$\flat$ & Mi$\flat$ & Sol$\flat$ & Mi$\flat$--Sol$\flat$--Do$\flat$ \\
Fa & Sol$\flat$ & Si$\flat$ & Re$\flat$ & Si$\flat$--Re$\flat$--Sol$\flat$ \\
\bottomrule
\end{tabularx}
\end{center}

\section{Regole di scrittura}
\label{sec:regole_napoletana}

\begin{regola}[Condotta della sesta napoletana]
	\label{reg:condotta_napoletana}
	\begin{itemize}
		\item Si \textbf{raddoppia la terza} (la nota al basso, IV grado). La fondamentale, cioè la nota napoletana, \textbf{non si raddoppia}.
		\item La nota napoletana \textbf{scende} alla sensibile per \textbf{terza diminuita} (Si$\flat$ $\searrow$ Sol$\sharp$ in La minore, Re$\flat$ $\searrow$ Si in Do maggiore). La terza diminuita è l'unico intervallo diminuito melodico ammesso in questo contesto. In alternativa la nota (o la terza raddoppiata) può essere \textbf{tenuta}, per diventare la settima della \sigla{V}{7}{}.
		\item La \textbf{quinta} dell'accordo (Fa in La minore, La$\flat$ in Do maggiore) scende di grado sulla fondamentale della dominante (Mi, Sol): si evita così la \textbf{seconda eccedente} Fa--Sol$\sharp$ (La$\flat$--Si in Do maggiore).
		\item Il \textbf{basso} sale di grado verso la dominante (Re $\nearrow$ Mi, Fa $\nearrow$ Sol).
		\item La terza raddoppiata scende (Re $\searrow$ Si, verso il V) oppure resta ferma come settima della \sigla{V}{7}{}.
	\end{itemize}
	Quando la terza raddoppiata scende sulla quinta del V (Re $\searrow$ Si) si forma una falsa relazione fra la nota napoletana e il II grado naturale (Si$\flat$ e Si in voci diverse). È l'unica falsa relazione tollerata dalla tradizione, perché la nota napoletana risolve sulla sensibile; si evita del tutto passando per il 6/4 di cadenza o tenendo la terza come settima della \sigla{V}{7}{}.
\end{regola}

\begin{esempio}[Sesta napoletana in La minore]
	\label{ex:napoletana_minore}
	A sinistra \sigla{II$\flat$}{6}{} $\rightarrow$ V, a destra \sigla{II$\flat$}{6}{} $\rightarrow$ \sigla{V}{7}{}: il Re del soprano (terza raddoppiata) è tenuto e diventa la settima della \sigla{V}{7}{}, che risulta incompleta (senza la quinta Si) e risolve sul Do. Il Si$\flat$ del tenore scende al Sol$\sharp$ per terza diminuita.
	\musica{es_cap21_02}
\end{esempio}

\begin{esempio}[Sesta napoletana e 6/4 di cadenza in La minore]
	\label{ex:napoletana_cadenza_minore}
	Nel soprano la nota napoletana disegna la linea La--Si$\flat$--La--Sol$\sharp$--La: il Si$\flat$ scende al La (nota del 6/4) e poi al Sol$\sharp$ della dominante. Il Fa scende al Mi, il Re al Do; il basso sale Re--Mi. La \sigla{V}{7}{} è completa e la tonica finale è incompleta.
	\musica{es_cap21_03}
\end{esempio}

\begin{esempio}[Sesta napoletana in Do maggiore]
	\label{ex:napoletana_maggiore}
	In tonalità maggiore la sesta napoletana richiede due alterazioni (Re$\flat$ e La$\flat$) che non sono in chiave: il Re$\flat$ scende a Si per terza diminuita, il La$\flat$ scende a Sol, il Fa al basso sale a Sol. Nella \sigla{V}{7}{} incompleta il Fa del contralto è tenuto.
	\musica{es_cap21_04}
\end{esempio}

\begin{esempio}[Sesta napoletana e 6/4 di cadenza in Do maggiore]
	\label{ex:napoletana_cadenza_maggiore}
	Nella \sigla{V}{7}{} completa (Sol-Si-Re-Fa) il Fa è la settima, che scende poi al Mi della tonica; il Sol del 6/4 resta fermo nel basso.
	\musica{es_cap21_05}
\end{esempio}

\begin{errore}[Raddoppiare la nota napoletana]
	\label{err:raddoppio_napoletana}
	Raddoppiare la fondamentale del II$\flat$ (Si$\flat$ in La minore) significa avere due note che devono entrambe scendere per terza diminuita sulla sensibile: si producono sensibile raddoppiata e ottave parallele.
\end{errore}

\begin{errore}[Seconda eccedente]
	\label{err:seconda_eccedente_napoletana}
	La quinta del \sigla{II$\flat$}{6}{} che sale alla sensibile (Fa $\nearrow$ Sol$\sharp$ in La minore, La$\flat$ $\nearrow$ Si in Do maggiore) produce un salto melodico di seconda eccedente, vietato.
\end{errore}

\begin{errore}[Sesta napoletana in stato fondamentale]
	\label{err:napoletana_fondamentale}
	La forma tipica e più usata è il primo rivolto. Lo stato fondamentale (II$\flat$) e il secondo rivolto compaiono solo in contesti più liberi e non si usano come soluzione abituale nel basso dato.
\end{errore}
